% Bộ đề mẫu Toán 11 chuẩn GDPT 2018 — Công văn 7991 & 1239
%[1D1N1-1]
\begin{ex}
Trên đường tròn lượng giác, góc lượng giác có số đo $\dfrac{\pi}{3}$ tương ứng với bao nhiêu độ?
\choice[4]
{\True $60^\circ$}
{$30^\circ$}
{$45^\circ$}
{$90^\circ$}
\loigiai{
Ta có $\dfrac{\pi}{3} = \dfrac{180^\circ}{3} = 60^\circ$. \textbf{Chọn A}.
}
\end{ex}

%[1D1H4-1]
\begin{ex}
Số nghiệm của phương trình $\sin x = \dfrac{1}{2}$ trên đoạn $[0;\pi]$ là
\choice[4]
{$1$}
{\True $2$}
{$3$}
{$0$}
\loigiai{
Trên $[0;\pi]$, $\sin x = \dfrac{1}{2} \Leftrightarrow x = \dfrac{\pi}{6}$ hoặc $x = \dfrac{5\pi}{6}$. Vậy có $2$ nghiệm. \textbf{Chọn B}.
}
\end{ex}

%[1D2N2-1]
\begin{ex}
Cho cấp số cộng $(u_n)$ có $u_1 = 3$ và công sai $d = 2$. Giá trị $u_5$ bằng
\choice[4]
{$9$}
{$10$}
{\True $11$}
{$13$}
\loigiai{
Ta có $u_5 = u_1 + 4d = 3 + 4 \cdot 2 = 11$. \textbf{Chọn C}.
}
\end{ex}

%[1D5H1-1]
\begin{ex}
Giá trị của giới hạn $\lim\limits_{n \to +\infty} \dfrac{2n+1}{n-3}$ bằng
\choice[4]
{$1$}
{$-\dfrac{1}{3}$}
{$0$}
{\True $2$}
\loigiai{
Ta có $\lim\limits_{n \to +\infty} \dfrac{2n+1}{n-3} = \lim\limits_{n \to +\infty} \dfrac{2 + \dfrac{1}{n}}{1 - \dfrac{3}{n}} = 2$. \textbf{Chọn D}.
}
\end{ex}

%[1D2H2-2]
\begin{ex}
Cho cấp số cộng $(u_n)$ có $u_1 = 2$ và $u_4 = 11$. Xét tính đúng sai của các mệnh đề sau:
\choiceTF[1t]
{\True Công sai của cấp số cộng là $d = 3$.}
{Số hạng thứ $10$ của cấp số cộng là $u_{10} = 30$.}
{\True Tổng $5$ số hạng đầu của cấp số cộng bằng $40$.}
{Số $100$ là một số hạng của cấp số cộng $(u_n)$.}
\loigiai{
\begin{itemchoice}
\itemch Đúng. Vì $u_4 = u_1 + 3d \Leftrightarrow 11 = 2 + 3d \Leftrightarrow d = 3$.
\itemch Sai. Vì $u_{10} = u_1 + 9d = 2 + 27 = 29$.
\itemch Đúng. Vì $S_5 = \dfrac{5(2u_1 + 4d)}{2} = \dfrac{5(4 + 12)}{2} = 40$.
\itemch Sai. Vì $2 + 3(n-1) = 100 \Leftrightarrow 3(n-1) = 98$ không có nghiệm nguyên dương $n$.
\end{itemchoice}
}
\end{ex}

%[1H4V1-2]
\begin{ex}
Cho hình chóp $S.ABCD$ có đáy là hình bình hành. Gọi $M$ là trung điểm $SC$. Giao điểm $I$ của đường thẳng $AM$ và mặt phẳng $(SBD)$ thỏa mãn $\overrightarrow{AI} = k \overrightarrow{AM}$. Tìm giá trị của $k$ (làm tròn đến hàng phần trăm).
\shortans[oly]{$0{,}67$}
\loigiai{
Gọi $O = AC \cap BD$. Trong mặt phẳng $(SAC)$, $I = AM \cap SO$ là trọng tâm tam giác $SAC$, suy ra $\overrightarrow{AI} = \dfrac{2}{3}\overrightarrow{AM}$, do đó $k = \dfrac{2}{3} \approx 0{,}67$.
}
\end{ex}

%[1D5V2-1]
\begin{bt}
Tính giới hạn $\lim\limits_{x \to 2} \dfrac{x^2 - 4}{x - 2}$.
\loigiai{
Ta có $\lim\limits_{x \to 2} \dfrac{x^2 - 4}{x - 2} = \lim\limits_{x \to 2} \dfrac{(x-2)(x+2)}{x-2} = \lim\limits_{x \to 2} (x+2) = 4$.
}
\end{bt}
